% Generated from locale/tr/content/second-order-logic/syntax-and-semantics/expressive-power.tex; do not hand-edit.


\olfileid[tr]{sol}{syn}{exp}
\olsection{İfade Gücü}

\begin{explain}
İkinci dereceden değişkenler üzerindeki niceleme, dilin ifade gücünü birinci
dereceden mantığın diline göre son derece artırır. İkinci dereceden varlık
nicelemesi, belirli özellikleri taşıyan fonksiyon ya da bağıntıların var
olduğunu söylememizi sağlar. Birinci dereceden mantıkta bunu yapmanın tek
yolu, bu amaç için mantıksal olmayan bir simge (yani !!a{function} ya da
bir !!{predicate}) belirlemektir. İkinci dereceden tümel niceleme, evrenin
bütün alt kümelerinin, onun üzerindeki bütün bağıntıların ya da !!{domain}
üzerinden yine !!{domain} üzerine bütün fonksiyonların bir
özelliği taşıdığını söylememizi sağlar. Birinci dereceden mantıkta yalnızca
dilin mantıksal olmayan simgelerinden birine atanan alt kümelerin,
bağıntıların ya da fonksiyonların bir özelliği taşıdığını söyleyebiliriz.
Üstelik bir özelliği taşıyan alt kümelerin, bağıntıların veya fonksiyonların
var olduğunu ya da bunların hepsinin o özelliği taşıdığını söylerken,
özelliğin kendisini belirlemek için de ikinci dereceden niceleme
kullanabiliriz. Böylece birinci dereceden mantıkta tanımlanamayan bağıntıları
tanımlayabilir ve evrenin birinci dereceden mantıkta ifade edilemeyen
özelliklerini ifade edebiliriz.
\end{explain}

\begin{defn}
$\Struct{M}$, bir~$\Lang{L}$ dili için !!a{structure} olsun.
$R \subseteq \Domain{M}^2$ bağıntısı, yalnızca $x$ ve~$y$ değişkenleri
serbest olan öyle bir !!{formula}~$!A_R(x,y)$ varsa~$\Lang{L}$ içinde
\emph{tanımlanabilir}dir ki $R(a,b)$'nin geçerli olması (yani
$\tuple{a,b} \in R$), $s(x)=a$ ve $s(y)=b$ koşulları altında ancak ve ancak
$\Sat{M}{!A_R(x,y)}[s]$ olduğunda gerçekleşsin.
\end{defn}

\begin{ex}
Birinci dereceden mantıkta özdeşlik bağıntısı~$\Id{\Domain{M}}$'yi
(yani $\Setabs{\tuple{a,a}}{a \in \Domain{M}}$ bağıntısını)
$\eq[x][y]$ formülüyle tanımlayabiliriz. İkinci dereceden mantıkta bu
bağıntıyı \emph{$\eq$ olmadan} tanımlayabiliriz. Çünkü $a$ ile $b$,
~$\Domain{M}$'nin aynı !!{element} öğesiyse~$\Domain{M}$'nin aynı alt
kümelerinin !!{element}s öğeleridir (zira kümeler !!{element}s öğeleriyle
belirlenir). Tersine, $a$ ile $b$ farklıysa aynı alt kümelerin !!{element}s
öğeleri değildir: örneğin $a \neq b$ ise $a \in \{a\}$, fakat
$b \notin \{a\}$ olur. Dolayısıyla ``~$\Domain{M}$'nin aynı alt
kümelerinin !!{element}s öğeleri olma'', $a$ ile $b$ için ancak ve ancak
$a=b$ olduğunda geçerli olan bir bağıntıdır. Bütün~$\Domain{M}$ alt kümeleri
üzerinde niceleme yapabildiğimiz için bu bağıntı ikinci dereceden mantıkta
ifade edilebilir. Öyleyse aşağıdaki !!{formula}, $\Id{\Domain{M}}$'yi
tanımlar:
\[
\lforall[X][(X(x) \liff X(y))]
\]
\end{ex}

\begin{prob}
$\lforall[X][(X(x) \lif X(y))]$ ifadesinin (dikkat: $\liff$ değil,
~$\lif$!) $\Id{\Domain{M}}$'yi tanımladığını gösterin.
\end{prob}

\begin{ex}
$R$, ikili bir !!{predicate} ise $\Assign{R}{M}$,~$\Domain{M}$ üzerinde
ikili bir bağıntıdır. Biraz kafa karıştırıcı olsa da hem !!{predicate}
olarak~$R$ için hem de bizzat $\Assign{R}{M}$ bağıntısı için $R$ harfini
kullanacağız. $R$'nin \emph{geçişli kapanışı}~$R^*$, bazı
$c_1$, \dots, $c_k$ için $R(a,c_1)$, $R(c_1,c_2)$, \dots, $R(c_k,b)$
olduğunda ve yalnızca o zaman $a$ ile $b$ arasında geçerli olan bağıntıdır.
Bu, $k=0$ durumunu da içerir; yani $R(a,b)$ ise $R^*(a,b)$ de olur.
Dolayısıyla $R \subseteq R^*$'dır. Aslında $R^*$,~$R$'yi içeren ve geçişli
olan en küçük bağıntıdır. İkinci dereceden mantıkta $X$'in~$R$'yi içeren
geçişli bir bağıntı olduğunu şöyle söyleyebiliriz:
\begin{multline*}
  !B_R(X) \ident \lforall[x][\lforall[y][(R(x,y) \lif X(x, y))]] \land {}\\
\lforall[x][\lforall[y][\lforall[z][((X(x,y) \land X(y,z)) \lif X(x,
      z))]]].
\end{multline*}
İlk bileşen $R \subseteq X$ olduğunu, ikincisi de $X$'in geçişli olduğunu
söyler.

$X$'in bu türden en küçük bağıntı olduğunu söylemek, kendisinin~$R$'yi
içeren ve geçişli olan her bağıntıda içerildiğini söylemektir. Böylece
~$R$'nin geçişli kapanışını şu !!{formula} ile tanımlayabiliriz:
\[
R^*(X) \ident !B_R(X) \land \lforall[Y][(!B_R(Y) \lif
  \lforall[x][\lforall[y][(X(x, y) \lif Y(x,y))]])].
\]
$\Sat{M}{R^*(X)}[s]$ ancak ve ancak $s(X)=R^*$ ise geçerlidir.
~$R$'nin geçişli kapanışı birinci dereceden mantıkta ifade edilemez.
\end{ex}

